\clearpage
\chapter*{Erratum to Exposé no. 195}
\addcontentsline{toc}{chapter}{Erratum to Exposé no. 195}
\setcounter{footnote}{0}
\src{300}

\noindent\textbf{4. TDTE II: The existence theorem in formal moduli theory.}

\begin{center}
[Bourbaki Seminar, vol. 12, 1959/60, no. 195, 22 pp.]
\end{center}

\noindent\textbf{Page 8.} --- The formula written in Proposition 5.1 is
correct only when $\Lambda$ is a field. In general,
$\mathfrak m_\xi/\mathfrak m_\xi^2$ must be replaced by the quotient of that
space by the image of $\mathfrak n_\xi/\mathfrak n_\xi^2$, where $\mathfrak n$
is the maximal ideal of $\Lambda$. Moreover, the definition given for $O$
smooth over $\Lambda$ is correct only when the residue-field extension $k'/k$
is separable. In the general case, cf. SGA III, 1.1.

\noindent\textbf{Page 14, remark.} --- The problems raised here are completely
solved in the projective case by “Hilbert schemes” (cf. TDTE IV). Examples of
Nagata and Hironaka show, however, that the functors under consideration need
not be representable without a projectivity hypothesis, even if one restricts
to classifying zero-dimensional subvarieties of a nonsingular complete
threefold. This is connected with the fact that the symmetric square of such a
variety may fail to exist.

\src{301}
\noindent\textbf{Page 15, no. 3.} --- For a fuller study, see TDTE V.

\noindent\textbf{Page 16, Proposition 3.1.} --- Read “if $f$ is proper and
separable” instead of “if $f$ is proper.”

\noindent\textbf{Pages 18 and 19, Remarks, 2.} --- I have recently proved that
the formal moduli scheme of an abelian variety over a field is indeed smooth
over the Witt ring; in other words, every abelian scheme over an Artinian local
ring that is a quotient of another ring is obtained by reduction from an
abelian scheme over the latter. The proof uses only the variance properties of
the obstruction class for lifting introduced in Exposé no. 182, page 12. Let
us also recall that the moduli schemes for curves of genus $g$ and for polarized
abelian schemes have been constructed by Mumford (cf. SHMT).

\noindent\textbf{Page 19, paragraph 5.} --- The section used to define $U$
must be assumed to be a non-zero-divisor not only on $\mathfrak X$, but on
every $X_n$.

\noindent\textbf{Page 20, lines 12 and 13.} --- Instead of
\[
  G_x=G\bigl(\operatorname{Spec}(\widehat{\mathcal O}_{\mathfrak X,x})\bigr),
\]
read
\[
  G_x=G\bigl(\operatorname{Spec}(\mathcal O_{\mathfrak X,x})\bigr)
  \subset
  G\bigl(\operatorname{Spec}(\widehat{\mathcal O}_{\mathfrak X,x})\bigr),
\]
and instead of “$\mathcal O'_{0x}$,” read
“$\widehat{\mathcal O}'_{0x}$.”
